\subsect{Source}

The dataset used in this assignment is the \hyperlink{https://www.kaggle.com/datasets/dgomonov/new-york-city-airbnb-open-data}{\textit{New York City Airbnb Open Data}} dataset, originally compiled by Dgomonov and publicly available on Kaggle.
The data was sourced from Inside Airbnb, an independent project that aggregates publicly available listing information from Airbnb's website. It reflects listing activity in New York City as of 2019.

\subsect{Size and Structure}

The dataset contains \textbf{48,895 rows} and \textbf{16 columns}. Each row represents a single Airbnb listing. The columns span a mix of categorical identifiers, geographic attributes, pricing and availability metrics, and host-level information.

\subsect{Key Columns}

Table~\ref{tab:columns} summarizes the columns most relevant to this assignment.

\begin{table}[h]
  \centering
  \begin{tabular}{lll}
    \hline
    \textbf{Column} & \textbf{Type} & \textbf{Description} \\
    \hline
    \texttt{neighbourhood\_group} & Categorical & Borough (Manhattan, Brooklyn, etc.) \\
    \texttt{neighbourhood}        & Categorical & Specific neighbourhood within borough \\
    \texttt{room\_type}           & Categorical & Entire home/apt, Private room, or Shared room \\
    \texttt{price}                & Numeric     & Nightly price in USD \\
    \texttt{minimum\_nights}      & Numeric     & Minimum required stay in nights \\
    \texttt{number\_of\_reviews}  & Numeric     & Total reviews received by listing \\
    \texttt{reviews\_per\_month}  & Numeric     & Average monthly review rate \\
    \texttt{availability\_365}    & Numeric     & Days available per year \\
    \texttt{calculated\_host\_listings\_count} & Numeric & Total listings by the same host \\
    \texttt{host\_name}           & Categorical & Name of the host \\
    \hline
  \end{tabular}
  \caption{Key columns in the NYC Airbnb dataset.}
  \label{tab:columns}
\end{table}

\subsect{Why This Dataset Requires Multi-Step Analysis}

No meaningful question about this dataset can be answered with a single operation.
The numeric columns -- such as \texttt{price}, \texttt{availability\_365}, and \texttt{reviews\_per\_month} -- are only interpretable in context: a high price means something different for a shared room in the Bronx than for an entire apartment in Manhattan. Extracting that context requires combining filtering, grouping, and aggregation across multiple categorical dimensions.

Furthermore, several quantities of analytical interest are not directly present as columns and must be derived before they can be used.
For example, estimated annual revenue potential must be computed as \texttt{price} \texorpdfstring{$\times$}{times} \texttt{availability\_365}, and minimum total cost requires multiplying \texttt{minimum\_nights} with \texttt{price}. These derived metrics must then feed into downstream grouping and ranking steps, making the pipeline inherently multi-stage.

The combination of geographic hierarchy (borough $\rightarrow$ neighbourhood), categorical room types, multiple numeric performance signals, and host-level aggregations creates a natural space for questions that span filtering, derivation, grouping, sorting, and top-$k$ selection -- all of which are required to produce a grounded, reproducible answer.
